\section{Honest intervals with a known score law}

A finite-label release determines how much of the assignment score is retained. We now consider interval estimation when its score law is known to the analyst. For trial size \leanref{sym:n}{\ensuremath{n}}, the \leanref{S-2}{known-score trial experiment} observes independent released rows under a chosen release. The sampling condition is stated explicitly.

\begin{assumptionv}{ass:iid-trial-sampling}[Independent trial sampling]
The trial observations are independent and identically distributed under \(P\):
\[
(R_i,A_i,Y_i)_{i=1}^n\stackrel{\mathrm{iid}}{\sim}\mathcal L_P(R,A,Y).
\]
\end{assumptionv}

\begin{definitionv}{synth_1}[Average treatment effect functional]
For any measure \(P\) on full rows \((E,Y_0,Y_1,A,Y,R)\) with score parameter \(\varepsilon\) and \(J\) labels, define the average treatment effect functional \(\theta(P)\) by
\[
\theta(P)=\int (Y_1-Y_0)\,dP.
\]
\end{definitionv}

Independent superpopulation sampling governs how the released-row law generates the trial data and is a standard condition for interval inference in partially identified models \citep{ImbensManski2004}.

Coverage must account for every causal law in \cref{obj:def:causal-law-class} consistent with the chosen release and score law. At miscoverage level \leanref{sym:alpha}{\ensuremath{\alpha}}, the following criterion allows the release and interval procedure to be selected jointly.

In the next two displays, \(P^{\otimes n}\) acts only on functions of the observed released rows. Its pushforward to those rows is \(\mathcal L_P(R,A,Y)^{\otimes n}\), the sampling law in \cref{obj:ass:iid-trial-sampling}.

\begin{definitionv}{def:honest-procedures}[Uniformly honest procedures $\mathcal J_{n,K,\alpha}(H)$]
\[
\mathcal J_{n,K,\alpha}(H)
=
\left\{
(g,C_n):
g\in\mathcal G_K,\quad
\inf_{P\in\mathfrak P(H,g)}
P^{\otimes n}\!\left[
\theta(P)\in C_n\bigl((R_i,A_i,Y_i)_{i=1}^n\bigr)
\right]
\ge 1-\alpha
\right\}.
\]
\end{definitionv}

For each pair in \leanref{sym:Jhonest}{\ensuremath{\mathcal J_{n,K,\alpha}(H)}}, the interval \leanref{sym:Cn}{\ensuremath{C_n}} attains the stated coverage uniformly over causal laws for its release. Among these pairs, the design criterion is worst-case expected interval length.

\begin{definitionv}{def:minimax-honest-length}[Minimax honest length $\mathcal L^\star_{n,K,\alpha}(H)$]
\[
\mathcal L^\star_{n,K,\alpha}(H)
=
\inf_{(g,C_n)\in\mathcal J_{n,K,\alpha}(H)}
\sup_{P\in\mathfrak P(H,g)}
\mathbb E_{P^{\otimes n}}\!\left[\operatorname{len}(C_n)\right].
\]
\end{definitionv}

The criterion \leanref{sym:Lminimax}{\ensuremath{\mathcal L^\star_{n,K,\alpha}(H)}} measures the shortest worst-case expected length attainable through joint choice of a release and an honest interval. The next result gives matching orders over the positive lower-density class in \cref{obj:def:score-lower-density-class}, together with a sampling lower bound for every score law.

\begin{theoremv}{thm:minimax-honest-length}[Minimax honest interval length]
Consider the following conditions:
\begin{itemize}
\item \textbf{(Overlap.)} The overlap parameter \(\varepsilon\) satisfies \cref{obj:ass:overlap}.
\item \textbf{(Miscoverage.)} \(0<\alpha<1/2\).
\item \textbf{(Trial size.)} \(n\geq1\).
\item \textbf{(Label budget.)} \(K\geq1\).
\end{itemize}
There are constants \(C_{\varepsilon,\alpha}>0\) and \(c_{\varepsilon,\alpha}>0\), independent of \(H,n,K\), such that every score law \(H\) on \(\mathcal E=[\varepsilon,1-\varepsilon]\) satisfies
\[
c_{\varepsilon,\alpha}n^{-1/2}
\leq \mathcal L^\star_{n,K,\alpha}(H),
\]
where the minimax expected length is defined in \cref{obj:def:minimax-honest-length}.

For every \(m_f>0\), there is a constant \(c_{\varepsilon,m_f,\alpha}>0\) such that every \(H\in\mathcal H^{\mathrm{lb}}_{\varepsilon,m_f}\), as defined in \cref{obj:def:score-lower-density-class}, satisfies
\[
c_{\varepsilon,m_f,\alpha}\bigl(K^{-1}+n^{-1/2}\bigr)
\leq \mathcal L^\star_{n,K,\alpha}(H)
\leq C_{\varepsilon,\alpha}\bigl(K^{-1}+n^{-1/2}\bigr).
\]

An interval attaining the upper bound uses the equal-width release with labels \(r\in\mathcal R_K=\{1,\ldots,K\}\):
\[
g_K(e)=1+\min\left\{\left\lfloor\frac{K(e-\varepsilon)}{1-2\varepsilon}\right\rfloor,K-1\right\},
\qquad
z_r=\varepsilon+\frac{(r-\tfrac12)(1-2\varepsilon)}{K}.
\]
For released trial rows \((R_i,A_i,Y_i)_{i=1}^n\), define
\[
\widehat T_n=\frac1n\sum_{i=1}^n
\left\{\frac{A_iY_i}{z_{R_i}}-\frac{(1-A_i)Y_i}{1-z_{R_i}}\right\},
\qquad
R_{n,K}=\frac4\varepsilon\sqrt{\frac{\log(4/\alpha)}{n}}
+\frac{1-2\varepsilon}{2K\varepsilon(1-\varepsilon)}.
\]
With \(\operatorname{clampATE}(t)=\max\{-1,\min\{1,t\}\}\), the interval
\[
C_n(x)=
\left[
\operatorname{clampATE}\bigl(\widehat T_n(x)-R_{n,K}\bigr),
\operatorname{clampATE}\bigl(\widehat T_n(x)+R_{n,K}\bigr)
\right]
\]
satisfies \((g_K,C_n)\in\mathcal J_{n,K,\alpha}(H)\), with honesty defined in \cref{obj:def:honest-procedures}. For every \(P\in\mathfrak P(H,g_K)\), as defined in \cref{obj:def:causal-law-class}, its expected length under the \(n\)-fold released-row law satisfies
\[
\int |C_n(x)|\,dP_{\mathrm{rel}}^{\otimes n}(x)
\leq C_{\varepsilon,\alpha}\bigl(K^{-1}+n^{-1/2}\bigr).
\]
\end{theoremv}

The two terms in \cref{obj:thm:minimax-honest-length} reflect score variation within a label cell and sampling variation across released rows. Midpoint inverse-probability weights give an honest interval whose radius accounts for both. Over the positive lower-density class, a label budget of order \(\sqrt n\) balances these terms and attains root-\(n\) expected length. The matching lower bound makes this balance a property of the experiment, rather than of the particular interval construction.
